\documentclass[11pt]{article}
\usepackage[utf8]{inputenc}
\usepackage{amsmath, amssymb, amsfonts}
\usepackage{geometry}
\usepackage{xcolor}
\geometry{a4paper, margin=1in}

\title{Physics-Aware Loss Function for VLSV-JAX Neural Correction}
\author{Vlasov-Jax Project Team}
\date{\today}

\begin{document}

\maketitle

\section{Introduction}
The objective of the neural corrector is to map a coarse-resolution ion distribution function $f_{coarse}(\mathbf{x}, \mathbf{v})$ to its high-fidelity counterpart $f_{fine}(\mathbf{x}, \mathbf{v})$ obtained from expensive fine-resolution simulations. To achieve this in a physically consistent manner, we utilize a multi-objective loss function that combines phase-space reconstruction accuracy with the preservation of macroscopic moments.

\section{Logarithmic Residual Mapping}
Given the high dynamic range of the distribution function (often spanning $10^{-6}$ to $0.5$), the MLP learns a logarithmic residual $\Delta \log f$:
\begin{equation}
    \Delta \log f(\mathbf{x}, \mathbf{v}) = \log\left( \frac{f_{fine}}{f_{coarse}} \right)
\end{equation}
The reconstructed distribution is then recovery via:
\begin{equation}
    \hat{f}_{fine} = f_{coarse} \cdot \exp\left( \text{MLP}(\mathbf{x}, \mathbf{v}; \theta) \right)
\end{equation}

\section{The Composite Loss Function}
The total loss $\mathcal{L}_{total}$ is defined as a weighted sum of the distribution loss and the physics-consistency loss:
\begin{equation}
    \mathcal{L}_{total} = \mathcal{L}_{dist} + \lambda_{phys} \mathcal{L}_{phys}
\end{equation}
where $\lambda_{phys}$ is a hyperparameter (typically set to 5.0) that governs the strength of the conservation constraints.

\subsection{Weighted Distribution Loss ($\mathcal{L}_{dist}$)}
To capture critical sub-grid kinetic phenomena such as reflected ion beams and high-energy tails, we employ a \textbf{Weighted Mean Squared Error (WMSE)}:
\begin{equation}
    \mathcal{L}_{dist} = \frac{1}{N_{samples}} \sum_{i} W(\mathbf{v}_i) \cdot \left\| (\Delta \log f)_{i}^{pred} - (\Delta \log f)_{i}^{target} \right\|^2
\end{equation}
The weighting kernel $W(\mathbf{v})$ emphasizes the high-velocity regions of phase-space:
\begin{equation}
    W(\mathbf{v}) = 1.0 + \frac{\|\mathbf{v}\|^2}{v_{scale}^2}
\end{equation}
This formulation ensures that the model does not ignore the low-density tails in favor of the high-density thermal core.

\subsection{Physics Consistency Loss ($\mathcal{L}_{phys}$)}
A critical requirement for any kinetic solver is the conservation of macroscopic quantities. We enforce \textbf{Mass and Momentum conservation} by penalizing the discrepancy between the moments calculated from the predicted distribution and those of the fine-resolution reference:
\begin{equation}
    \mathcal{L}_{phys} = \|n[\hat{f}] - n[f_{fine}]\|^2 + \|\mathbf{V}[\hat{f}] - \mathbf{V}[f_{fine}]\|^2
\end{equation}
where the moments are calculated via numerical integration over the velocity grid:
\begin{align}
    n[f] &= \int f(\mathbf{v}) \, d^3\mathbf{v} \\
    \mathbf{V}[f] &= \frac{1}{n} \int \mathbf{v} f(\mathbf{v}) \, d^3\mathbf{v}
\end{align}
In our JAX implementation, these integrals are computed efficiently using vectorized tensor operations, ensuring that the conservation laws are differentiable and can guide the optimization of the network weights.

\end{document}
